\begin{table}[htbp]
\centering
\caption{问题2.3：最优方案站点实际有效负荷、附件容量与弹性校验}
\begin{tabular}{|c|c|c|c|c|c|c|c|}
\hline
站点 & 规模 & 附件容量 & 弹性上限 & 实际有效负荷 & 附件利用率 & S2 & 容量状态 \\
\hline
A & 大型 & 3000 & 3060 & 3049.0 & 101.6\% & 0.60 & 弹性内超载1.63\% \\
C & 小型 & 1000 & 1020 & 1015.4 & 101.5\% & 0.60 & 弹性内超载1.54\% \\
E & 大型 & 3000 & 3060 & 2980.7 & 99.4\% & 0.60 & 未超附件容量 \\
\hline
全区 & — & 7000 & 7140 & 7045.1 & 100.6\% & — & 弹性上限内 \\
\hline
\end{tabular}
\vspace{0.2cm}
说明：方案B中站点负荷为归属小区的实际有效日服务人次之和；本轮按附件容量的 2\% 小幅弹性上限校验。
\end{table}
